\section{引言}

部分子分布函数（PDF）是理解强子结构以及对高能散射实验的截面作出预言的关键量。在硬散射过程的 QCD 因子化定理中~\cite{Collins:1989gx}，相关的 PDF 定义为光锥关联算符在核子态之间的矩阵元。例如，在维数正规化 $d=4-2\epsilon$ 下，裸的非极化夸克 PDF 为
\begin{align} \label{eq:lcpdf}
	q(x,\epsilon) \equiv\! \int\! {d\xi^-\over 4\pi} e^{-ixP^+\xi^-}\! \langle P | \bar{\psi}  (\xi^-) \gamma^+ U (\xi^-, 0) \psi (0) | P \rangle ,
\end{align}
其中 $x$ 是动量分数，核子动量 $P^\mu=(P^0,0,0,P^z)$，$\xi^{\pm}=(t\pm z)/\sqrt{2}$ 是光锥坐标，而 Wilson 线为
\begin{align}
	U(\xi^-,0) = P\exp\biggl(-ig\int_0^{\xi^-} d\eta^- A^+(\eta^-)\biggr)\ .
\end{align}
通常把裸 PDF 在 $\overline{\ensuremath{\operatorname{MS}}}$ 方案下重正化得到 $q(x,\mu)$，再用重正化后的 PDF 对实验作预言。其关系为
\beq
	q(x,\epsilon) = \int_x^1 {dy\over y}\:  Z^{\overline{\ensuremath{\operatorname{MS}}}}\Bigl({x\over y},\epsilon,\mu\Bigr) q(y,\mu)\ ,
\eeq
其中 $\mu$ 是重正化标度，我们在重正化常数 $Z^{\overline{\ensuremath{\operatorname{MS}}}}$ 与 PDF 中略去了味道指标。在 $A^+=0$ 的光锥量子化中，$\overline{\ensuremath{\operatorname{MS}}}$ 定义可以解释为部分子数密度。

迄今为止，关于 PDF 的主要知识来自对深度非弹性散射与喷注数据的全局拟合，例如参见文献~\cite{Ball:2014uwa,Dulat:2015mca,Martin:2009iq,Alekhin:2017kpj,Buckley:2014ana}。另一方面，用 QCD 从第一性原理计算 PDF 一直是一个有吸引力的课题，例如它可以提供难以通过实验确定的自旋与动量分布信息。为此人们考虑了几种基于格点理论的方法——格点理论是求解 QCD 的一种非微扰方法。由于格点理论定义在具有虚时间方向的离散欧氏空间上，要计算像 PDF 这类依赖实时的闵可夫斯基量非常困难。第一个也是研究得最充分的方法是计算 PDF 的矩（moments）~\cite{Martinelli:1987zd,Martinelli:1988xs,Detmold:2001dv,Detmold:2002nf,Dolgov:2002zm}，它们是定域规范不变算符的矩阵元。然而，由于格点正规化破坏了 $O(4)$ 转动对称性，不同维数算符之间的混合使高阶矩难以计算，实践中限制了该方法所能提取的信息量。文献~\cite{Davoudi:2012ya}提出了一种通过恢复转动对称性来改善这一状况的方法。
其他建议包括从强子张量~\cite{Liu:1993cv,Liu:1998um,Liu:1999ak,Liu:2016djw}、前向 Compton 振幅~\cite{Chambers:2017dov}提取 PDF，可能借助味改变流~\cite{Detmold:2005gg}，以及更一般的``格点截面''~\cite{Ma:2014jla,Ma:2017pxb}。对这些方法的系统性格点分析仍在进行中，但挑战犹存。

文献~{\cite{Ji:2013dva}} 中 Ji 提出，PDF 的 $x$ 依赖性可以从格点上的一个欧氏分布中提取，这可以用大动量有效理论（LaMET）的语言来理解~\cite{Ji:2014gla}。该欧氏分布被称为准 PDF（quasi-PDF），其裸矩阵元由沿 $z$ 方向的夸克空间关联定义：
\beq \label{eq:qpdf}
	\tilde{q}(x,P^z,\epsilon) \equiv \int_{-\infty}^\infty {dz\over 4\pi} e^{ixP^zz}\langle P | \bar{\psi}  (z) \Gamma U (z, 0) \psi (0) | P \rangle \,,
\eeq
其中 $\Gamma=\gamma^z$，$z^\mu=ze^\mu$，$e^\mu=(0,0,0,1)$，Wilson 线为
\beq
	U(z,0) = P\exp\left(-ig\int_0^z dz' A^z(z')\right)\ .
\eeq
对有限动量 $P^z$，$\tilde{q}(x,P^z,\epsilon)$ 的支撑区为 $-\infty < x < \infty$。
按照文献~\cite{Hatta:2013gta}，可以考虑一类具有普适性的算符。例如对准 PDF，也可以把式~(\ref{eq:qpdf}) 中的 $\Gamma=\gamma^z$ 换成 $\Gamma=\gamma^0$，因为两种定义在沿 $z$ 方向的无穷洛伦兹助推下都约化为 PDF。与式~(\ref{eq:lcpdf})中在洛伦兹助推下不变的 PDF 不同，准 PDF 通过核子动量 $P^z$ 动态地依赖于助推。当核子动量 $P^z$ 远大于核子质量 $M$ 与 $\Lambda_{\rm QCD}$ 时——这在格点上是一个可以达到的窗口——准 PDF 可以因子化为一个匹配系数与 PDF 的乘积~\cite{Ji:2013dva,Ji:2014gla}。因子化公式为
\begin{align} \label{eq:momfact}
	\tilde{q}(x,P^z,\mu_R)
	=& \int_{-1}^1 {dy\over |y|}\: C\Bigl({x\over y},{\mu_R\over \mu},{\mu\over p^z}\Bigr) \, q(y,\mu)
	\\
	& + \mathcal{O}\biggl({M^2\over P_z^2},{\Lambda_{\ensuremath{\operatorname{QCD}}}^2\over P_z^2}\biggr)\ ,\nn
\end{align}
其中重正化准 PDF $\tilde{q}(x,P^z,\mu_R)$ 在某一特定方案、重正化标度 $\mu_R$ 下定义，而 $\mathcal{O}(M^2/P_z^2, \Lambda_{\rm QCD}^2/P_z^2)$ 是被核子动量压低的幂次修正。一般地，$C$ 的结果会依赖于 $\Gamma=\gamma^z$ 或 $\gamma^0$ 的选择以及重正化方案。对 $\Gamma=\gamma^z$，匹配系数 $C$ 已在单圈层次对同位旋矢量夸克准 PDF 作出计算：最初带横向动量截断~\cite{Xiong:2013bka}，随后在文献~\cite{Ma:2014jla,Alexandrou:2015rja}中得到确认，最近又在正规化不变动量减除（RI/MOM）方案中被确定~\cite{Stewart:2017tvs}。胶子准 PDF 的匹配在文献~\cite{Wang:2017qyg,Wang:2017eel}中计算。匹配系数 $C$ 与用于 $\tilde q$ 和 $q$ 的态的选择无关。\footnote{方案的定义本身可能单独涉及态的选择，例如 RI/MOM 方案，但其结果仍是算符重正化，可以用于 $\tilde q$ 和 $q$ 的不同强子态选择。横向截断方案与 $\overline{\rm MS}$ 方案甚至不需要这种选择。}由于匹配计算是用动量为 $p^z$ 的夸克态进行的，$C$ 应当怎样取值可能并不显然，一些文献建议在对强子核子态利用 $C$ 时取 $p^z=P^z$。例如最初的准 PDF 论文~\cite{Ji:2013dva,Ji:2014gla,Xiong:2013bka}以及从准 PDF 得到 PDF 的开创性格点计算~\cite{Lin:2014zya,Alexandrou:2015rja,Chen:2016utp,Alexandrou:2016jqi,Zhang:2017bzy,Alexandrou:2017huk,Chen:2017mzz,Lin:2017ani,Chen:2017gck}就是如此，这些在文献~\cite{Lin:2017snn}中有总结。在准广义部分子分布的分析~\cite{Ji:2015qla}中曾观察到应取 $p^z=|y|P^z$。通过对式~(\ref{eq:momfact}) 的严格分析，我们将证明该方程的正确结果确实是 $p^z=|y|P^z$。

最近，一种不同的做法~\cite{Radyushkin:2017cyf}被提出：从与文献~\cite{Ji:2013dva}相同的格点 QCD 矩阵元中提取 PDF，但基于空间关联函数（或称赝 PDF 方法）的洛伦兹不变变量，而不是准 PDF。在这一方法中，人们从对 $\mu=0$ 或 $\mu=z$ 定义的\spcorr $\tilde{Q}_{\gamma^\mu}$ 出发：
\begin{align}
	{1\over2}\langle P | \tilde{O}_{\gamma^\mu}(z,\epsilon) | P\rangle
	= & P^\mu \tilde{Q}_{\gamma^\mu}(\zeta=P^zz,z^2,\epsilon) \,,
\end{align}
它依赖于两个洛伦兹不变量 $z^2$ 与 $\zeta=-z\cdot P = P^z z$；后者也称为 Ioffe 时间。对任意狄拉克矩阵 $\Gamma$，算符 $\tilde{O}_\Gamma$ 定义为
\begin{align} \label{eq:OGamma}
	\tilde{O}_{\Gamma}(z,\epsilon)  = \bar{\psi}  (z) \Gamma U (z, 0) \psi (0) \,.
\end{align}
这正是用于定义准 PDF 的同一空间关联函数（可在格点上计算）（见式~(\ref{eq:qpdf})），那里固定 $P^z$ 并对 $z$ 作傅里叶变换。如果改为固定 $z^2$，并把 Ioffe 时间 $\zeta$——原则上是对 $P^z$ 积分——傅里叶变换到动量分数 $x$，就得到赝 PDF（pseudo-PDF）~\cite{Radyushkin:2017cyf}，
\begin{align} \label{eq:pseudo}
\mathcal{P} \left( x, z^2,\epsilon\right)
   = \int_{-\infty}^\infty \! \frac{d \zeta}{2\pi}\:
    e^{ix \zeta}\: \tilde{Q}_{\gamma^0}\bigl( \zeta, z^2,\epsilon \bigr)
   \,.
\end{align}
对任意有限的 $z$，赝 PDF 只有 $-1\le x \le 1$ 的支撑区~\cite{Radyushkin:1983wh,Radyushkin:2016hsy}，但没有部分子模型解释。
（同样，赝 PDF 也完全可以取 $\Gamma=\gamma^z$。）\spcorr/赝 PDF 方法已在格点上被探索~\cite{Orginos:2017kos,Karpie:2017bzm}，其中研究了短距离行为。
PDF 对应于 $z^\mu$ 为类光的情形，此时时空关联函数被称为 Ioffe 时间分布~\cite{Ioffe:1969kf}，
\begin{align}
 Q(\zeta,\epsilon) = \tilde Q_{\gamma^+}(\zeta=-P^+\xi^-,z^2=0,\epsilon)
  \,.
\end{align}
经傅里叶变换后，该关联给出 PDF
\begin{align} \label{eq:Qdefn}
q(y,\epsilon) =\int_{-\infty}^{\infty} {d\zeta\over 2\pi} \ e^{iy\zeta}\: Q(\zeta,\epsilon)\,.
\end{align}
简言之，准 PDF 和赝 PDF 是欧氏空间关联函数的不同表示，总结于表~\ref{tab:I}。
\begin{table}
\begin{center}
	\begin{tabular}{ | c | c | c |}
		\hline
		分布 & 自\spcorr 的傅里叶变换 & 宗量\\[-2pt]
		&  & \\ \hline
		空间关联  &  & $\zeta = zP^z, z^2$\\ \hline
		准 PDF & $z\to xP^z$ & $x, P^z$\\ \hline
		赝 PDF & $\zeta \to x$ & $x, z^2$\\ \hline
	\end{tabular}
\caption{各欧氏分布之间关系的总结。}
\label{tab:I}
\end{center}
\end{table}

文献~\cite{Ji:2017rah}指出，为了从小 $z^2$ 的\spcorr 获得足够的信息来提取 PDF，必须进行大动量 $P^z$ 的格点计算，这与准 PDF 的要求相同。文献~\cite{Ji:2017rah}还提出，重正化赝 PDF 满足以下小 $z^2$ 因子化，
\begin{align} \label{eq:pseudofact}
	\mathcal{P} \left( x, z^2\mu_R^2\right)
  &= \int {dy\over |y|}\: \mathcal{C}\left({x\over y}, {\mu_R^2\over \mu^2},z^2\mu_R^2\right) q(y,\mu)
   \\
  &\ + {\cal O}(z^2\Lambda_{\rm QCD}^2, z^2 M^2) \,,\nn
\end{align}
并且他们在 $\Gamma=\gamma^0$ 的非极化同位旋矢量情形以 $O(\alpha_s)$ 精度验证了它。（同样，系数 ${\cal C}$ 会依赖于 $\Gamma=\gamma^0$ 或 $\gamma^z$ 的选择。）

文献~\cite{Ma:2014jla}给出了准 PDF 因子化公式即式~(\ref{eq:momfact})的一个图解推导。本文我们以另一种方式推导准 PDF 的这一因子化公式，并证明\spcorr 和赝 PDF 是同一基本因子化的不同表示。我们的方法基于类空间隔定域算符的算符乘积展开（OPE）~\cite{Wilson:1969zs}。对这类算符，OPE 在纯量场论中已被证明到微扰论的所有阶~\cite{Brandt:1967rb,Wilson:1972ee,Zimmermann:1972tv}，并被普遍假定对包括 QCD 在内的任何可重正化量子场论成立。通过引入辅助场代替 Wilson 线~\cite{Dorn:1986dt}，已知式~(\ref{eq:OGamma})中的关联函数等价于此类型定域可重正化算符的乘积。通过我们的推导，我们得到了大 $P^z$ 与小 $z^2$ 因子化公式——式~(\ref{eq:momfact})与式~(\ref{eq:pseudofact})——的显式形式，以及匹配系数 $C$ 与 $\mathcal{C}$ 之间的关系。由于大 $P^z$ 与小 $z^2$ 的要求对准 PDF 和赝 PDF 两种方法是相同的，原则上用任何一种方法来做 $\tilde{O}_{\Gamma}(z)$ 质子矩阵元的格点计算都没有根本差别。用相同的格点数据比较两种方法是有趣的，尽管它们原则上不会给出不同的结果。

本文其余部分安排如下：在第~\ref{sec:ope}~节中，我们用 $\tilde{O}_\Gamma(z)$ 的 OPE 推导式~(\ref{eq:momfact})中 LaMET 的大 $P^z$ 因子化和式~(\ref{eq:pseudofact})中赝 PDF 的小 $z^2$ 因子化。我们证明在式~(\ref{eq:momfact})中必须取 $p^z=|y|P^z$，因此 $C$ 中相应的宗量是 $\mu/(|y|P^z)$。
（这种 OPE 方法最近在文献~\cite{Ma:2017pxb}中被用于证明``格点截面''的因子化定理；这里完成的 OPE 证明是独立完成的，并首次发表于文献~\cite{Zhao:2017int}。）
%
在第~\ref{sec:equiv}~节中，我们在 $\overline{\rm MS}$ 下单圈推导\spcorr、赝 PDF 与准 PDF 分布及其匹配系数，并分析准 PDF 与赝 PDF 之间的傅里叶变换关系。与早先 $\overline{\rm MS}$ 下准 PDF 的结果不同，我们还使用最小减除的维数正规化来重正化 $x=\pm\infty$ 处的发散。第~\ref{sec:ren}~节讨论如何把 $\overline{\rm MS}$ 以外的重正化方案方便地纳入因子化公式。
%
第~\ref{sec:num}~节我们对匹配系数作数值分析：用全局拟合确定的 PDF~\cite{Martin:2009iq}数值计算式~(\ref{eq:momfact})中的卷积。我们证明，在式~(\ref{eq:momfact})中使用 $p^z=P^z$ 与 $p^z=|y|P^z$ 的差别是一个重要的效应，而且我们的 $\overline{\rm MS}$ 匹配系数对卷积积分中的截断并不敏感。第~\ref{sec:lattice}~节讨论我们的 OPE 分析对准 PDF 与赝 PDF 两种方法中 PDF 格点计算的影响。最后，我们在第~\ref{sec:summary}~节给出结论。
